\documentclass[11pt]{article}
\usepackage[margin=1in]{geometry}
\usepackage{amsmath,amssymb,amsthm,mathtools,booktabs,array,enumitem}
\usepackage[T1]{fontenc}
\usepackage{lmodern}
\usepackage{microtype}
\usepackage[hypertexnames=false]{hyperref}
\hypersetup{colorlinks=true,linkcolor=blue,citecolor=blue,urlcolor=blue,
 pdfauthor={Christopher D. Long},
 pdftitle={A Critical Linear Bound with Logarithmic Savings for Glazer's Time-Limited Fair-Coin Problem}}
\setlength{\emergencystretch}{2em}

\newtheorem{theorem}{Theorem}[section]
\newtheorem{proposition}[theorem]{Proposition}
\newtheorem{lemma}[theorem]{Lemma}
\newtheorem{corollary}[theorem]{Corollary}
\theoremstyle{definition}
\newtheorem{definition}[theorem]{Definition}
\newtheorem{example}[theorem]{Example}
\theoremstyle{remark}
\newtheorem{remark}[theorem]{Remark}

\newcommand{\N}{\mathbb N}
\newcommand{\Z}{\mathbb Z}
\newcommand{\Ftwo}{\mathbb F_2}
\newcommand{\eps}{\varepsilon}
\newcommand{\wt}{\operatorname{wt}}
\newcommand{\sgn}{\operatorname{sgn}}
\newcommand{\poly}{\mathbb R[z]_{\le n}}
\newcommand{\E}{\mathbb E}
\newcommand{\Prb}{\mathbb P}
\newcommand{\logfive}{\log_5}

\title{A Critical Linear Bound with Logarithmic Savings\\
for Glazer's Time-Limited Fair-Coin Problem}
\author{Christopher D. Long\\[3pt]
\small\href{mailto:galizur@gmail.com}{\texttt{galizur@gmail.com}}}
\date{September 20, 2026}

\begin{document}
\maketitle

\begin{abstract}
Let a coin have unknown bias $p\in(0,1)$, and let $L$ be the length of its
maximal initial constant run. We construct a deterministic, bias-independent
fair-bit extractor satisfying, simultaneously,
\[
 T(L)\le\lceil C_*L\rceil,
 \qquad
 T(L)\le C_*L-\log_5L+6,
 \qquad
 C_*:=1+\frac{2\log\phi}{\log5}=1.597987435665\ldots,
\]
where $\phi=(1+\sqrt5)/2$.
The construction separately balances dyadic blocks at horizon $2N$.
A parity-character folding identity supplies a real center; an elementary
binomial-mixture estimate controls its positive-level norms; and a
boundary-preserving integer rounding moves each interior coefficient by at
most five. The argument is uniform for dyadic $N\ge128$, and six smaller
blocks have exact integer certificates.
We prove that $C_*$ is the optimal leading coefficient for all separately
balanced horizon-$2N$ rules, without assuming complement symmetry, and
that this lower bound persists with an $o(N)$ extra horizon. An exact
inequality for the deadlines bounds the possible logarithmic saving.
The extractor also works for exchangeable binary sources with almost surely
finite $L$. Finally, a parity obstruction characterizes why the block
length must be a power of two at horizon $2N$.
\end{abstract}

\begingroup
\small
\makeatletter
\renewcommand{\l@section}[2]{\@dottedtocline{1}{0em}{1.8em}{#1}{#2}}
\makeatother
\tableofcontents
\endgroup
\clearpage

\section{Statement of the problem and main result}

Let $X_1,X_2,\ldots$ be independent coin flips with
\[
 \mathbb P(X_j=0)=p,
 \qquad
 \mathbb P(X_j=1)=q:=1-p,
\]
where $p\in(0,1)$ is arbitrary and unknown to the rule.  On every nonconstant sequence, let $L$ denote the length of the maximal initial constant run.  Thus the observed sequence starts with either $0^L1$ or $1^L0$.

A deterministic extractor assigns one of two outputs, denoted $H$ and $T$, after reading finitely many flips.  It is \emph{fair} if
\[
 \mathbb P_p(H)=\mathbb P_p(T)=\frac12
 \qquad\text{for every }p\in(0,1).
\]
Let $T(L)$ be an upper bound for the number of flips used on sequences whose initial run has length $L$.

Put
\begin{equation}\label{eq:constants}
 \phi:=\frac{1+\sqrt5}{2},
 \qquad
 \beta:=\frac{2\log\phi}{\log5},
 \qquad
 \alpha:=1-\beta,
 \qquad
 C_*:=1+\beta=2-\alpha.
\end{equation}
Numerically,
\[
 \beta=0.597987435665\ldots,
 \qquad
 \alpha=0.402012564334\ldots,
 \qquad
 C_*=1.597987435665\ldots.
\]

For each power of two $N\ge2$, define the integer saving
\begin{equation}\label{eq:shift}
 s_N:=\max\{0,\lfloor\log_5N\rfloor-3\}.
\end{equation}
The logarithm $\log_5 x$ means $\log x/\log5$.

\begin{theorem}[Critical bound with logarithmic saving]\label{thm:main}
There exists a deterministic fair extractor that stops on every nonconstant
input sequence. It uses two flips when $L=1$. When $N\le L<2N$ with $N\ge2$
a power of two, it satisfies
\begin{equation}\label{eq:main-profile}
 T(L)\le\min\{2N,\lceil C_*L\rceil-s_N\}.
\end{equation}
In particular, for every integer $L\ge1$, the same extractor satisfies
\begin{align}
 T(L)&\le\lceil C_*L\rceil,\label{eq:main-bound}\\
 T(L)&\le C_*L-\log_5L+6.\label{eq:main-log-bound}
\end{align}
Its two outputs have equal mass separately on every dyadic block for every
bias $p\in(0,1)$.
\end{theorem}

The classical assignments $01\mapsto H$ and $10\mapsto T$ treat $L=1$;
they go back to von Neumann~\cite{vonNeumann1951}. A substantial literature
studies more efficient extraction from a source of unknown bias, typically
measuring expected source cost, block efficiency, or asymptotic entropy rate;
see, for example, Hoeffding--Simons~\cite{HoeffdingSimons1970},
Elias~\cite{Elias1972}, Stout--Warren~\cite{StoutWarren1984},
Peres~\cite{Peres1992}, Juels--Jakobsson--Shriver--Hillyer
\cite{JuelsEtAl2000}, and Pae--Loui~\cite{PaeLoui2005}. The efficiency
criterion here is different: it is a pathwise worst-case deadline indexed by
the observed initial-run length, and is required uniformly over the unknown
bias.

The particular time-limited problem considered here was posed by Elliot
Glazer in an online discussion in August 2024, initially for a loaded
six-sided die, together with the question of the best linear efficiency
constant~\cite{GlazerPuzzling2024}. There are two incompatible historical
conventions for the parameter called $n$. In the Puzzling discussion and in
Wagon's answer, let us write $n_{\mathrm{SE}}$ for the number of flips up to
and including the first outcome different from the initial one. Thus
\begin{equation}\label{eq:historical-notation}
 n_{\mathrm{SE}}=L+1.
\end{equation}
For the binary case, Charles Wang gave a corrected
$2n_{\mathrm{SE}}-2=2L$ construction~\cite{WangPuzzling2024}. Glazer
pointed out a flaw in the initially posted version; Wang agreed and edited
the construction to make the recursion greedy. That corrected construction
is the valid one cited here. Wang also proposed a recursive $7n_{\mathrm{SE}}/4$
refinement, but that refinement was subsequently found to be invalid;
Wagon's later account records that Peter Winkler pointed out the error and
that Wang acknowledged it~\cite{WagonPuzzling2025}. Stan Wagon subsequently
reported using the problem as a Macalester Problem of the Week; several
readers found a recursive solution, and Wagon and Peter Winkler found a
formulaic construction. Wagon's posted account gives the bound
$2n_{\mathrm{SE}}-3=2L-1$ for the binary problem~\cite{WagonPuzzling2025}.

Glazer later posed the coin version as Problem 12592 in the \emph{American
Mathematical Monthly}~\cite{GlazerAMM2026}. There the symbol $n$ has a
different meaning: the ``critical value'' is the length of the maximal
initial constant run. Writing this historical parameter as
$n_{\mathrm{AMM}}$, we have
\[
 n_{\mathrm{AMM}}=L=n_{\mathrm{SE}}-1.
\]
Consequently the two earlier bounds become
\[
 2n_{\mathrm{SE}}-2=2L=2n_{\mathrm{AMM}},
 \qquad
 2n_{\mathrm{SE}}-3=2L-1=2n_{\mathrm{AMM}}-1.
\]
Thus the Monthly requests $2n_{\mathrm{AMM}}$ and
$\max\{2,2n_{\mathrm{AMM}}-1\}$ are exactly the same two bounds after
translating conventions. This paper uses $L$ throughout.

For $L\ge2$ we balance the events $N\le L<2N$ individually. This is a
sufficient construction principle, not an assumption on all possible fair
extractors. Section~\ref{sec:optimality} proves optimality of the leading
coefficient within the separately balanced horizon-$2N$ class, and
Section~\ref{sec:extra-horizon} treats sublinear additional horizons.
Neither result asserts unrestricted optimality when different blocks can
cancel each other's imbalance.

The distinction between an observed stopping time and a worst-case profile
is important: $T(L)$ bounds every continuation with that initial run length.
A universally fair extractor cannot stop at a constant prefix $0^t$ or
$1^t$. Indeed, an output assigned to $0^t$ would have probability at least
$p^t$, exceeding $1/2$ for $p$ sufficiently near one; the other case is
symmetric. Thus observing the first change, at flip $L+1$, is unavoidable.

\section{The dyadic packet reduction}

Fix a dyadic integer $N\ge2$, put
\[
 n:=N-1,
\]
and consider initial-run lengths $N+i$, where $0\le i\le n$.  Pad all decisions in this block to the common horizon $2N$.  After the first change, the full tail then has length $n-i$.

For a tail $w\in\{0,1\}^{n-i}$, pair
\[
 0^{N+i}1w
 \qquad\text{with}\qquad
 1^{N+i}0\overline w,
\]
where $\overline w$ is the bitwise complement.  If $w$ has $r$ zeros, then a positively signed pair has probability difference
\[
 p^{N+i+r}q^{N-i-r}-q^{N+i+r}p^{N-i-r}.
\]
Writing $\ell=i+r$, this depends only on $\ell$.  The class $\ell=0$ contributes zero, and fairness is obtained by giving every class $1\le\ell\le n$ zero signed discrepancy.

Suppose that at level $i$ the final $a_i$ tail bits are left unread.  Put
\[
 R_i:=n-i-a_i.
\]
Thus $R_i$ tail bits are read, and the actual stopping length at this
level is
\[
 (N+i)+1+R_i=2N-a_i.
\]
For each read prefix $x\in\{0,1\}^{R_i}$ choose a sign
$\sigma_i(x)\in\{\pm1\}$.  The same sign is used for all $2^{a_i}$ unread completions, while the complementary branch receives the opposite output.  For $0\le r\le R_i$, let $e_{i,r}$ be the number of positive read prefixes containing $r$ zeros minus the number of negative such prefixes, and write
\[
 E_i(x):=\sum_{r=0}^{R_i}e_{i,r}x^r.
\]

\begin{lemma}[Packet realization]\label{lem:packet-realization}
An integer vector $(e_{i,r})_{r=0}^{R_i}$ is realizable by deterministic signs if and only if
\begin{equation}\label{eq:packet-admissibility-e}
 |e_{i,r}|\le\binom{R_i}{r},
 \qquad
 e_{i,r}\equiv\binom{R_i}{r}\pmod2
 \qquad(0\le r\le R_i).
\end{equation}
The signed discrepancy contributed by level $i$ to class $\ell$ is
\[
 [x^\ell]x^i(1+x)^{a_i}E_i(x).
\]
Consequently, the block is fair whenever
\begin{equation}\label{eq:x-identity}
 \sum_{i=0}^{n}x^i(1+x)^{a_i}E_i(x)
\end{equation}
is constant.
\end{lemma}

\begin{proof}
There are $Q=\binom{R_i}{r}$ read prefixes with $r$ zeros.  If $P$ of them are positive, their net coefficient is $2P-Q$.  This proves necessity of \eqref{eq:packet-admissibility-e}; conversely, those two conditions make $(Q+e_{i,r})/2$ an integer in $[0,Q]$.  Declaring, for example, the lexicographically first $(Q+e_{i,r})/2$ prefixes positive gives a deterministic realization.

A read prefix with $r$ zeros has $\binom{a_i}{j}$ unread completions with $j$ zeros.  Such a completion belongs to class $i+r+j$.  Summing over $r$ and $j$ gives the coefficient formula.  If \eqref{eq:x-identity} is constant, all nonzero classes have discrepancy zero.  The constant class is harmless because its paired probability difference is identically zero.
\end{proof}

Substitute $x=-z$ and define
\[
 f_{i,r}:=(-1)^{i+r}e_{i,r}.
\]
The capacity and parity conditions are unchanged.  It is therefore enough to find integers $f_{i,r}$ such that
\begin{equation}\label{eq:target-identity}
 \sum_{i=0}^{n}\sum_{r=0}^{R_i}
 f_{i,r}z^{i+r}(1-z)^{a_i}=1
\end{equation}
and
\begin{equation}\label{eq:packet-admissibility}
 |f_{i,r}|\le\binom{R_i}{r},
 \qquad
 f_{i,r}\equiv\binom{R_i}{r}\pmod2.
\end{equation}

\begin{lemma}[Fixed-horizon type balance]\label{lem:type-balance}
Let $\mathcal W\subseteq\{0,1\}^{H}$ and let
$\varepsilon:\mathcal W\to\{\pm1\}$ be a signing. The identity
\[
 \sum_{w\in\mathcal W}\varepsilon(w)p^{z(w)}(1-p)^{H-z(w)}=0
 \qquad(0<p<1)
\]
holds if and only if
\begin{equation}\label{eq:type-balance}
 \sum_{\substack{w\in\mathcal W\\z(w)=k}}\varepsilon(w)=0
 \qquad(0\le k\le H),
\end{equation}
where $z(w)$ is the number of zeros in $w$.
\end{lemma}
\begin{proof}
Write the latter signed counts as $c_k$. Dividing the former expression
by $(1-p)^H$ gives $\sum_k c_k x^k=0$ for every
$x=p/(1-p)>0$. A nonzero real polynomial has only finitely many roots,
so every $c_k$ is zero. The converse follows by grouping the words by $k$.
\end{proof}

\section{The logarithmically shifted suffix profile}\label{sec:profile}

We first record elementary facts about the constants. The number $\beta$
is irrational: if $\beta=A/B$ with positive integers $A,B$, then
$5^A=\phi^{2B}$, whereas
\[
 \phi^{2B}=F_{2B}\phi+F_{2B-1}\notin\mathbb Q.
\]
Here $F_0=0$, $F_1=1$, and $F_{j+1}=F_j+F_{j-1}$; the identity follows
inductively from $\phi^2=\phi+1$. In particular, both $\alpha$ and $C_*$
are irrational.

We shall use
\begin{equation}\label{eq:alpha-bounds}
 \frac25<\alpha<\frac{13}{32}.
\end{equation}
For the first inequality it suffices to show $\phi^{10}<5^3$.
Since $\phi<13/8$,
\[
 \phi^{10}=55\phi+34<55\cdot\frac{13}{8}+34=\frac{987}{8}<125.
\]
For the second, it suffices to show $\phi^{64}>5^{19}$. The recurrence gives
$F_{64}=10610209857723$ and $F_{63}=6557470319842$, while
$\phi>809/500$ follows from $(559/250)^2<5$. Hence
\[
 500\phi^{64}>809F_{64}+500F_{63}
 =11862394934818907>9536743164062500=500\cdot5^{19}.
\]

For the rest of the construction use $s=s_N$ from \eqref{eq:shift} and set
\begin{equation}\label{eq:critical-profile}
 a_i:=\max\{0,2N-\lceil C_*(N+i)\rceil+s\},
 \qquad R_i:=n-i-a_i.
\end{equation}
Equivalently,
\begin{equation}\label{eq:profile-floor}
 a_i=\max\{0,\lfloor\alpha(N+i)\rfloor-2i+s\}.
\end{equation}
In particular the suffixes are specified by integers and the sign rule
does not depend on the unknown bias.

\begin{lemma}[Saving and profile properties]\label{lem:profile}
Let $N\ge2$ be a power of two and let $h=\min\{i:a_i=0\}$. Then:
\begin{enumerate}[label=\textup{(\roman*)}]
\item $s_N=0$ for $N<128$. For $N\ge128$,
\begin{equation}\label{eq:shift-bounds}
 5^{s_N}\le\frac N{125},\qquad 0\le s_N\le\frac N{64}.
\end{equation}
\item $0\le a_i\le n-i$ and $a_i-a_{i+1}\in\{0,1,2\}$.
\item For $0\le i<h$,
\begin{equation}\label{eq:R-beta}
 R_i=\lfloor\beta(N+i)\rfloor-s_N,
\end{equation}
and $R_0,\ldots,R_h$ is nondecreasing.
\item The stopping length is exactly
\begin{equation}\label{eq:profile-stop}
 2N-a_i=\min\{2N,\lceil C_*(N+i)\rceil-s_N\}.
\end{equation}
\item For dyadic $N\ge128$,
\begin{equation}\label{eq:kR-bounds}
 k:=a_0\le\frac{27N}{64}-1,\qquad
 R_0=n-k\ge\frac{37N}{64}.
\end{equation}
\item For $N\ge8$, $h\le N/2$, $a_{n-2}=a_{n-1}=a_n=0$, and
$R_i\ge2$ for $0\le i\le n-2$.
\end{enumerate}
\end{lemma}
\begin{proof}
Put $j=\lfloor\log_5N\rfloor$. When $128\le N<625$, $j=3$ and $s=0$,
so the first inequality in \eqref{eq:shift-bounds} holds. When $N\ge625$,
$s=j-3$ and $5^s=5^j/125\le N/125$. The elementary induction
$64(j-3)\le5^j$ for $j\ge4$ proves the second inequality; for $s=0$
it is immediate. The assertion for $N<128$ follows from $N<5^4$.

The untruncated expression in \eqref{eq:profile-floor} drops by either one
or two at each step. Truncation at zero gives (ii)'s drop assertion.
If $a_i>0$, irrationality and $\alpha+\beta=1$ give
\[
 R_i=N+i-1-\lfloor\alpha(N+i)\rfloor-s
     =\lfloor\beta(N+i)\rfloor-s.
\]
For $N<128$ this is nonnegative because $s=0$. For $N\ge128$,
$\beta>19/32$ and $s\le N/64$ give
$R_i>37N/64-1>0$. Thus $a_i\le n-i$ in every case, including the
zero-suffix levels. The displayed $R_i$ are nondecreasing before the
transition. At the transition,
\[
 R_h-R_{h-1}=a_{h-1}-1\in\{0,1\}.
\]
This proves (ii)--(iii). Assertion (iv) is the definition of $a_i$.

For dyadic $N\ge128$, $13N/32$ is integral, so
\[
 k=\lfloor\alpha N\rfloor+s
 \le\frac{13N}{32}-1+\frac N{64}
 =\frac{27N}{64}-1.
\]
This proves (v).

For $N\ge128$, the untruncated expression at $i=N/2$ is less than
$39N/64-N+N/64<0$. For $8\le N<128$, the same assertion follows with
$s=0$. Hence $h\le N/2\le n-2$. Before $h$, the read lengths are at
least $R_0$; for $8\le N<128$ we have $R_0=\lfloor\beta N\rfloor\ge4$,
and for larger $N$ use (v). After $h$, $R_i=n-i$, proving (vi).
\end{proof}

\section{A parity-character folding identity}

The central construction is most transparent after the substitution
\begin{equation}\label{eq:y-u}
 y:=\frac{x}{1+x},
 \qquad
 u:=1-2y=\frac{1-x}{1+x}.
\end{equation}
A parity character on $t$ specified read coordinates has expectation $u^t$ under independent Bernoulli-$y$ inputs.  Equivalently,
\begin{equation}\label{eq:parity-character}
 (1+x)^{R}u^t=(1-x)^t(1+x)^{R-t}.
\end{equation}
Every coefficient of the polynomial on the right has absolute value at most the corresponding coefficient of $(1+x)^R$.

For a polynomial $G(u)=\sum_{t\ge0}p_tu^t$ with $p_t\ge0$, define its fold at degree $R$ by
\begin{equation}\label{eq:fold}
 \mathcal F_R G(u)
 :=\sum_{t=0}^{R-1}p_tu^t+
   \left(\sum_{t\ge R}p_t\right)u^R.
\end{equation}

Set
\begin{equation}\label{eq:F0}
 F_0(u):=\left(\frac{1+u}{2}\right)^n.
\end{equation}
We construct polynomials $W_i,F_{i+1}$ recursively.  Suppose
\[
 F_i(u)=\eps_i\sum_t p_{i,t}u^t,
 \qquad
 \eps_i\in\{\pm1\},
 \qquad
 p_{i,t}\ge0.
\]
If $\deg F_i\le R_i$, put $W_i:=F_i$ and $W_j:=0$ for $j>i$.  Otherwise put
\begin{equation}\label{eq:fold-recursion}
 W_i:=\eps_i\mathcal F_{R_i}\!\left(\sum_tp_{i,t}u^t\right),
 \qquad
 F_{i+1}:=\frac{F_i-W_i}{y}
          =\frac{2(F_i-W_i)}{1-u}.
\end{equation}

\begin{lemma}[One-sign recursion]\label{lem:one-sign}
If $\deg F_i>R_i$, then
\begin{equation}\label{eq:Fnext-coeff}
 F_{i+1}(u)
 =-\eps_i\sum_{s\ge R_i}
   \left(2\sum_{t>s}p_{i,t}\right)u^s.
\end{equation}
Thus the coefficients of $F_{i+1}$ again have one sign, and
\[
 \deg F_{i+1}=\deg F_i-1.
\]
The recursion terminates exactly at level $h$, and
\begin{equation}\label{eq:fold-telescope}
 \left(\frac{1+u}{2}\right)^n
 =\sum_{i=0}^{h}y^iW_i(u).
\end{equation}
\end{lemma}

\begin{proof}
For $G(u)=\sum_tp_tu^t$,
\begin{align*}
 G(u)-\mathcal F_RG(u)
 &=\sum_{t>R}p_t(u^t-u^R)\\
 &=-(1-u)\sum_{s\ge R}
      \left(\sum_{t>s}p_t\right)u^s.
\end{align*}
This proves \eqref{eq:Fnext-coeff}.  Initially $\deg F_0=n$, and every nonterminal division lowers the degree by one.  Hence $\deg F_i=n-i$.  By definition, $R_i=n-i-a_i$, so the first level at which $\deg F_i\le R_i$ is exactly the first level with $a_i=0$, namely $h$.  Finally, $F_i=W_i+yF_{i+1}$; iteration gives \eqref{eq:fold-telescope}.
\end{proof}

Write
\[
 W_i(u)=\sum_{t=0}^{R_i}w_{i,t}u^t,
 \qquad
 M_i:=\sum_{t=0}^{R_i}|w_{i,t}|.
\]
Define a real packet polynomial
\begin{equation}\label{eq:real-E}
 E_i^{(0)}(x)
 :=(1+x)^{R_i}W_i\!\left(\frac{1-x}{1+x}\right)
 =\sum_{t=0}^{R_i}w_{i,t}(1-x)^t(1+x)^{R_i-t}.
\end{equation}
After replacing $x$ by $-z$, write
\[
 E_i^{(0)}(-z)=(-1)^i\sum_{r=0}^{R_i}u_{i,r}z^r.
\]

\begin{proposition}[Exact real center]\label{prop:real-center}
The coefficients $u_{i,r}$ satisfy
\begin{equation}\label{eq:real-identity}
 \sum_{i=0}^{n}\sum_{r=0}^{R_i}
 u_{i,r}z^{i+r}(1-z)^{a_i}=1.
\end{equation}
Moreover,
\begin{equation}\label{eq:center-capacity-M}
 |u_{i,r}|\le M_i\binom{R_i}{r},
\end{equation}
$u_{0,0}=1$, and $u_{i,r}=0$ for all $i>h$.
\end{proposition}

\begin{proof}
Since $(1+u)/2=1-y$, equation \eqref{eq:fold-telescope} reads
\[
 (1-y)^n=\sum_i y^iW_i(u).
\]
Substitute $y=x/(1+x)$ and multiply by $(1+x)^n$.  Because $n-i-R_i=a_i$, the result is
\[
 1=\sum_i x^i(1+x)^{a_i}E_i^{(0)}(x).
\]
Replacing $x$ by $-z$ gives \eqref{eq:real-identity}.

For fixed $t$, every coefficient of $(1-x)^t(1+x)^{R_i-t}$ is a signed sum of exactly $\binom{R_i}{r}$ terms of magnitude one.  The triangle inequality gives \eqref{eq:center-capacity-M}.  Finally, $u_{0,0}=W_0(1)=F_0(1)=1$, and termination gives the last assertion.
\end{proof}

\section{Elementary binomial estimates and logarithmic slack}\label{sec:norms}

All expectations in this section are finite binomial sums. The estimate
below uses neither Fourier inversion nor a concentration theorem.

Let $X\sim\operatorname{Bin}(n,1/2)$.  For $i\ge1$, define
\[
 \binom{m_+}{i}:=
 \begin{cases}
 \binom mi,&m\ge i,\\
 0,&m<i.
 \end{cases}
\]

\begin{lemma}[Exact folded mass]\label{lem:mass-formula}
For $1\le i\le h$,
\begin{equation}\label{eq:mass-formula}
 M_i=2^i\,\mathbb E
 \binom{(X-R_{i-1})_+}{i}.
\end{equation}
Also $M_0=1$.
\end{lemma}

\begin{proof}
Let $p_{0,x}=2^{-n}\binom nx$.  Repeated use of \eqref{eq:Fnext-coeff}, together with the monotonicity of $R_0,\ldots,R_h$, gives
\begin{equation}\label{eq:iterated-tail}
 p_{i,s}
 =2^i\sum_{x>s}p_{0,x}\binom{x-s-1}{i-1}
 \qquad(s\ge R_{i-1}).
\end{equation}
This follows by induction: summing \eqref{eq:iterated-tail} over $t>s$ uses
\[
 \sum_{t=s+1}^{x-1}\binom{x-t-1}{i-1}
 =\binom{x-s-1}{i}.
\]
Summing \eqref{eq:iterated-tail} over $s\ge R_{i-1}$ and applying the same hockey-stick identity gives
\[
 M_i
 =2^i\sum_xp_{0,x}\binom{x-R_{i-1}}i,
\]
which is \eqref{eq:mass-formula}.  The equality $M_0=1$ follows directly from \eqref{eq:F0} and the fold preserving the sum of absolute coefficients.
\end{proof}

\begin{lemma}[An elementary binomial atom bound]\label{lem:atom-bound}
Let $Z\sim\operatorname{Bin}(n,\eta)$, where $0<\eta<1$, and put
$\rho=2\min\{\eta,1-\eta\}$. Then
\begin{equation}\label{eq:atom-bound}
 \max_k\Prb(Z=k)^2\le\frac1{(n+1)\rho}.
\end{equation}
\end{lemma}
\begin{proof}
Write $c_m=2^{-m}\binom{m}{\lfloor m/2\rfloor}$ for the largest atom
of $\operatorname{Bin}(m,1/2)$. The usual consecutive-coefficient ratio
locates this maximum. We have
\begin{equation}\label{eq:fair-atom-square}
 c_m^2\le\frac1{m+1}.
\end{equation}
Indeed, $c_0=1$ and
$c_{2r+2}=\frac{2r+1}{2r+2}c_{2r}$; induction for even indices uses
$(2r+1)(2r+3)<(2r+2)^2$. For odd indices,
$c_{2r+1}=\frac{2r+1}{2r+2}c_{2r}$, and
$(2r+1)/(2r+2)^2\le1/(2r+2)$.

By symmetry it suffices to consider $\eta\ge1/2$, so
$\rho=2(1-\eta)$. Generate each trial by using a fair Bernoulli trial
with probability $\rho$ and the deterministic value one otherwise.
If $M\sim\operatorname{Bin}(n,\rho)$ counts the fair trials, then
conditionally on $M=m$ the sum has distribution
$n-m+\operatorname{Bin}(m,1/2)$. This mixture identity also follows by
expanding
\[
 (1-\eta+\eta z)^n
 =\left((1-\rho)z+\rho\frac{1+z}{2}\right)^n.
\]
Consequently, finite Cauchy--Schwarz and \eqref{eq:fair-atom-square} give
\[
 \max_k\Prb(Z=k)^2
 \le\bigl(\E c_M\bigr)^2
 \le\E c_M^2\le\E\frac1{M+1}.
\]
Finally, $\binom nm/(m+1)=\binom{n+1}{m+1}/(n+1)$ gives the exact identity
\[
 \E\frac1{M+1}
 =\frac{1-(1-\rho)^{n+1}}{(n+1)\rho}
 \le\frac1{(n+1)\rho}.
\]
The formula remains valid when $\rho=1$. This proves the lemma.
\end{proof}

Put
\begin{equation}\label{eq:tilting-constants}
 p_*:=1-\frac1{\sqrt5},\qquad q_*:=\frac1{\sqrt5},\qquad
 \eta:=\frac{\sqrt5}{1+\sqrt5}.
\end{equation}
We use the identities
\begin{equation}\label{eq:critical-identities}
 \frac2{\sqrt5-1}=\frac{1+\sqrt5}{2}=\phi,
 \qquad (\sqrt5)^\beta=\phi.
\end{equation}

\begin{lemma}[Tilted mass identity]\label{lem:tilted-mass}
For $i\ge1$ define, on integers,
\[
 b_i(m):=
 \begin{cases}
 \binom mi p_*^i q_*^{m-i},&m\ge i,\\
 0,&m<i.
 \end{cases}
\]
If $Z\sim\operatorname{Bin}(n,\eta)$ and $1\le i\le h$, then
\begin{equation}\label{eq:tilted-mass}
 M_i=(\sqrt5)^{\beta(n+i)-R_{i-1}}
          \E b_i(Z-R_{i-1}).
\end{equation}
Moreover, for every finite integer $K\ge i$,
\begin{equation}\label{eq:kernel-mass}
 \sum_{m=i}^{K}b_i(m)\le\frac1{p_*}.
\end{equation}
\end{lemma}
\begin{proof}
For $R=R_{i-1}$ and $x-R\ge i$,
\[
 2^i\binom{x-R}{i}
 =\phi^i(\sqrt5)^{x-R}b_i(x-R).
\]
Also exponential tilting of $X\sim\operatorname{Bin}(n,1/2)$ gives,
for every function $g$ on its finite support,
\[
 \E(\sqrt5)^Xg(X)=\phi^n\E g(Z).
\]
Substituting these identities into \eqref{eq:mass-formula} and using
$(\sqrt5)^\beta=\phi$ proves \eqref{eq:tilted-mass}.

For \eqref{eq:kernel-mass}, the number
$p_*b_i(m)=\binom mi p_*^{i+1}q_*^{m-i}$ is the probability that the
$(i+1)$st success occurs at trial $m+1$ in $K+1$ Bernoulli trials.
These events are disjoint as $m$ ranges from $i$ to $K$, so their
probabilities sum to at most one. Equivalently, this is the finite identity
\[
 p_*\sum_{m=i}^{K}b_i(m)
 =\sum_{j=i+1}^{K+1}\binom{K+1}{j}p_*^j q_*^{K+1-j}\le1.
\]
No infinite-series argument is required.
\end{proof}

\begin{proposition}[Uniform half-norm bound]\label{prop:digit-norms}
For every dyadic $N\ge128$,
\begin{equation}\label{eq:half-norm}
 M_0=1,\qquad M_i<\frac12\quad(1\le i\le h).
\end{equation}
\end{proposition}
\begin{proof}
For $1\le i\le h$, the preceding level has positive suffix, so
\[
 R_{i-1}=\lfloor\beta(n+i)\rfloor-s_N,
 \qquad \beta(n+i)-R_{i-1}<s_N+1.
\]
By \eqref{eq:kernel-mass},
\[
 \E b_i(Z-R_{i-1})\le\frac1{p_*}\max_k\Prb(Z=k).
\]
Apply Lemma~\ref{lem:atom-bound} with
$\rho=2(1-\eta)=2/(1+\sqrt5)$ and $n+1=N$ to obtain
\begin{equation}\label{eq:mass-square-bound}
 M_i^2<\frac{5(1+\sqrt5)}{2p_*^2}\frac{5^{s_N}}N
       <27\frac{5^{s_N}}N\le\frac{27}{125}<\frac14.
\end{equation}
Every numerical comparison here is elementary: $\sqrt5<9/4$ and
$p_*>11/20$ imply
\[
 \frac{5(1+\sqrt5)}{2p_*^2}<\frac{3250}{121}<27,
\]
and $5^{s_N}\le N/125$ is \eqref{eq:shift-bounds}.
The equality $M_0=1$ was proved in Lemma~\ref{lem:mass-formula}.
\end{proof}

\section{Boundary-preserving integral parity completion}

We now prove an abstract rounding lemma.  For packet sites write
\[
 v_{i,r}:=z^{i+r}(1-z)^{a_i}.
\]
Put
\[
 d_i:=a_i-a_{i+1}\in\{0,1,2\}.
\]
For $0\le i<n$ and $1\le r\le R_i$, define the integral relation
\begin{equation}\label{eq:kernel-relation}
 K_{i,r}
 :=\mathbf e_{i,r}
 -\sum_{j=0}^{d_i}(-1)^j\binom{d_i}{j}
  \mathbf e_{i+1,r-1+j}.
\end{equation}
All indicated child sites exist because $R_{i+1}=R_i+d_i-1$.  The represented polynomial is zero.

\begin{lemma}[Integral kernel basis]\label{lem:kernel-basis}
The relations $K_{i,r}$ form a basis, over $\mathbb R$, of the kernel of
\[
 A:(c_{i,r})\longmapsto\sum_{i,r}c_{i,r}v_{i,r}
 \in\mathbb R[z]_{\le n}.
\]
Moreover, there exists an integral vector $q=(q_{i,r})$ such that
\begin{equation}\label{eq:q-identity}
 Aq=1,
 \qquad
 q_{i,r}\equiv\binom{R_i}{r}\pmod2.
\end{equation}
\end{lemma}

\begin{proof}
The boundary columns
\[
 v_{i,0}=z^i(1-z)^{a_i},
 \qquad0\le i\le n,
\]
form a unitriangular basis of $\mathbb R[z]_{\le n}$.  Thus $A$ has rank $n+1$.  The number of relations \eqref{eq:kernel-relation} is
\[
 \sum_{i=0}^{n-1}R_i
 =\sum_{i=0}^{n}(R_i+1)-(n+1),
\]
the dimension of the kernel.  They are linearly independent because $K_{i,r}$ has coefficient one at its parent site $(i,r)$ and all remaining coordinates lie at level $i+1$.

For the second assertion, start with the integer vector
\[
 c_{i,r}:=\binom{R_i}{r}.
\]
Modulo two, its polynomial image is
\begin{align*}
 Ac
 &\equiv\sum_{i=0}^{n}z^i(1+z)^{a_i+R_i}\\
 &=\sum_{i=0}^{n}z^i(1+z)^{n-i}\\
 &=(1+z)^{n+1}-z^{n+1}\\
 &\equiv1\pmod2,
\end{align*}
because $n+1=N$ is a power of two.  Hence $Ac=1+2H$ for some $H\in\mathbb Z[z]_{\le n}$.  The unitriangular boundary basis expresses $H$ as an integral combination of the $v_{i,0}$.  Subtract twice that combination from $c$ to obtain $q$.
\end{proof}

The following lemma is the integral step that permits the critical coefficient.  Its exact treatment of the two packet boundaries is essential.

\begin{lemma}[Five-unit lattice rounding]\label{lem:lattice-rounding}
Assume that $N\ge8$ is a power of two. Let an integer profile satisfy
$0\le a_i\le n-i$, $a_i-a_{i+1}\in\{0,1,2\}$,
$a_{n-2}=a_{n-1}=a_n=0$, and $R_i=n-i-a_i\ge2$ for $i\le n-2$.
Let $u=(u_{i,r})$ be a real solution of \eqref{eq:real-identity} such that
\begin{equation}\label{eq:u-box}
 |u_{i,r}|\le\binom{R_i}{r},
 \qquad
 u_{0,0}=1,
\end{equation}
and $u_{i,r}=0$ on the final two levels $i=n-1,n$.  Then there is an integral solution $f$ of \eqref{eq:target-identity} with the required packet parities such that:
\begin{enumerate}[label=\textup{(\roman*)}]
\item every lower and upper packet boundary is exact:
\begin{equation}\label{eq:boundary-exact}
 f_{i,0},f_{i,R_i}\in\{\pm1\};
\end{equation}
\item at every nonboundary site,
\begin{equation}\label{eq:five-error}
 |f_{i,r}-u_{i,r}|\le5;
\end{equation}
\item if $a_{i-1}=a_i=0$ and $u_{i,r}=0$, then at every nonboundary site of level $i$,
\begin{equation}\label{eq:tail-error}
 |f_{i,r}|\le2.
\end{equation}
\end{enumerate}
\end{lemma}

\begin{proof}
Choose $q$ as in Lemma~\ref{lem:kernel-basis}.  Since $A(q-u)=0$, there are unique real coefficients $\gamma_{i,r}$ such that
\[
 q-u=\sum_{i=0}^{n-1}\sum_{r=1}^{R_i}\gamma_{i,r}K_{i,r}.
\]
We shall choose even integers $m_{i,r}$ and put
\[
 \delta_{i,r}:=\gamma_{i,r}-m_{i,r},
 \qquad
 f:=q-\sum_{i,r}m_{i,r}K_{i,r}
   =u+\sum_{i,r}\delta_{i,r}K_{i,r}.
\]
Then $Af=1$, and $f$ has the same packet parities as $q$.

For every ordinary kernel coordinate, choose $m_{i,r}$ to be a nearest even integer to $\gamma_{i,r}$, so that
\begin{equation}\label{eq:delta-bound}
 |\delta_{i,r}|\le1.
\end{equation}
We modify this choice along the two boundary chains, always preserving \eqref{eq:delta-bound}.

For the lower boundary, observe that
\[
 f_{i+1,0}=u_{i+1,0}-\delta_{i,1}.
\]
Choose $b_{i+1}\in\{\pm1\}$ nearest to $u_{i+1,0}$ and prescribe
\begin{equation}\label{eq:lower-delta}
 \delta_{i,1}:=u_{i+1,0}-b_{i+1}.
\end{equation}
This has absolute value at most one.  It belongs to the permitted coset $\gamma_{i,1}+2\mathbb Z$: indeed, the lower-boundary coordinate of $q-u$ gives
\[
 \gamma_{i,1}=u_{i+1,0}-q_{i+1,0},
\]
and $q_{i+1,0}-b_{i+1}$ is even.  We make this prescription for $0\le i\le n-2$; the final overlapping variable is treated below.

For the upper boundary, write $\eta_i:=\delta_{i,R_i}$ and process $0\le i\le n-2$ in increasing order.  At level zero choose $b_0^+\in\{\pm1\}$ so that
\[
 \eta_0=b_0^+-u_{0,R_0}
\]
has absolute value at most one.  At a later level choose $b_i^+\in\{\pm1\}$ so that
\begin{equation}\label{eq:upper-delta}
 \eta_i=b_i^+-u_{i,R_i}+(-1)^{d_{i-1}}\eta_{i-1}
\end{equation}
has absolute value at most one.  This is possible because the expression from which $b_i^+$ is subtracted lies in $[-2,2]$.  It also lies in the permitted coset.  Indeed, the upper-boundary coordinate of $q-u$ is
\[
 q_{i,R_i}-u_{i,R_i}
 =\gamma_{i,R_i}-(-1)^{d_{i-1}}\gamma_{i-1,R_{i-1}},
\]
so the difference between \eqref{eq:upper-delta} and $\gamma_{i,R_i}$ is even by induction.  This makes $f_{i,R_i}=b_i^+$.

By the structural assumptions, $R_{n-2}=2$, $R_{n-1}=1$, and $R_n=0$,
while $R_i\ge2$ for $i\le n-2$. The only overlap of the
lower and upper kernel variables is therefore $(i,r)=(n-1,1)$.  The
preceding upper parameter is $\eta_{n-2}=\delta_{n-2,R_{n-2}}$, where
$R_{n-2}=2$.  Because the real center vanishes on the final two levels,
the upper coordinate at level $n-1$ and the lower coordinate at level $n$
give
\[
 q_{n-1,1}+q_{n,0}=-\gamma_{n-2,R_{n-2}}.
\]
The left side is even.  Hence $\gamma_{n-2,R_{n-2}}$ is an even integer, and so is $\eta_{n-2}$; together with $|\eta_{n-2}|\le1$, this forces
\[
 \eta_{n-2}=0.
\]
Moreover, the lower coordinate at level $n$ gives
$q_{n,0}=-\gamma_{n-1,1}$; since $q_{n,0}$ is odd,
$\gamma_{n-1,1}$ is an odd integer.  We may therefore choose
$\delta_{n-1,1}=1$ or $-1$ in its allowed odd coset.  The two final
boundary values are then opposite signs.  This completes
\eqref{eq:boundary-exact} and preserves \eqref{eq:delta-bound}
everywhere.

At an arbitrary site, the error $f-u$ consists of at most one parent term of magnitude one and the child image of the preceding level, whose total coefficient magnitude is at most
\[
 \sum_{j=0}^{d_{i-1}}\binom{d_{i-1}}j\le4.
\]
This proves \eqref{eq:five-error}.  If both consecutive suffix lengths are zero, then $d_{i-1}=0$ and the center vanishes; the error is a difference of at most two numbers of absolute value one, proving \eqref{eq:tail-error}.
\end{proof}

\section{Capacity reserve at level zero}

For levels $i\ge1$, Proposition~\ref{prop:digit-norms} already gives proportional reserve.  The level-zero digit has norm exactly one, so it requires a separate estimate.

Let
\[
 R:=R_0=n-k,
 \qquad
 X\sim\operatorname{Bin}(n,1/2),
 \qquad
 T:=\min\{X,R\}.
\]
The fold $W_0$ gives the distribution
\[
 \mathbb P(T=t)=2^{-n}\binom nt\quad(0\le t<R),
 \qquad
 \mathbb P(T=R)=\mathbb P(X\ge R).
\]
Let
\begin{equation}\label{eq:krawtchouk}
 K_t^{(R)}(r):=[x^r](1-x)^t(1+x)^{R-t}.
\end{equation}
The substitution $x=-z$ contributes the factor $(-1)^r$ to the
coefficient of degree $r$.  Hence
\begin{equation}\label{eq:u0-kraw}
 u_{0,r}=(-1)^r\mathbb E K_T^{(R)}(r),
 \qquad
 |u_{0,r}|=\left|\mathbb E K_T^{(R)}(r)\right|.
\end{equation}

\begin{lemma}[Swap estimate]\label{lem:swap}
For $1\le t\le R-1$ and $1\le r\le R-1$,
\begin{equation}\label{eq:swap-bound}
 \left|K_t^{(R)}(r)\right|
 \le\binom Rr-2\binom{R-2}{r-1}.
\end{equation}
\end{lemma}

\begin{proof}
Fix a $t$-subset $A$ of an $R$-element set.  Then
\[
 K_t^{(R)}(r)=\sum_{|B|=r}(-1)^{|A\cap B|}.
\]
Choose $a\in A$ and $b\notin A$.  On the family of $r$-subsets containing exactly one of $a,b$, the involution swapping $a$ and $b$ reverses the sign, so those terms cancel.  There are $2\binom{R-2}{r-1}$ such subsets.  Bounding the remaining terms by their number proves the claim.
\end{proof}

\begin{lemma}[Level-zero reserve]\label{lem:level-zero-reserve}
For every dyadic $N\ge128$ and every $1\le r\le R_0-1$,
\begin{equation}\label{eq:level-zero-slack}
 \binom{R_0}{r}-|u_{0,r}|>5.
\end{equation}
\end{lemma}
\begin{proof}
Set $R=R_0$, $k=a_0$, and $X\sim\operatorname{Bin}(n,1/2)$ as above.
The elementary binomial moment identities give
\[
 \E X=\frac n2,\qquad \E(X-n/2)^2=\frac n4.
\]
They follow, for example, by summing
$x\binom nx=n\binom{n-1}{x-1}$ and
$x(x-1)\binom nx=n(n-1)\binom{n-2}{x-2}$.
For
\[
 \mathcal A=\left\{\frac n3\le X\le\frac{n-1}{2}\right\},
\]
symmetry for odd $n$ gives $\Prb(X\le(n-1)/2)=1/2$. On $X<n/3$,
$(X-n/2)^2>n^2/36$, whence
\[
 \Prb(X<n/3)\le\frac{9}{n}<\frac14\qquad(n\ge127).
\]
Therefore
\begin{equation}\label{eq:A-prob}
 \Prb(\mathcal A)>\frac14.
\end{equation}
On this event, the truncated variable $T=\min\{X,R\}$ equals $X$ and
lies between $1$ and $R-1$.

For $2\le r\le R-2$, the swap estimate gives
\[
 \binom Rr-|u_{0,r}|
 \ge2\Prb(1\le T\le R-1)\binom{R-2}{r-1}
 >\frac12\binom{R-2}{r-1}\ge\frac{R-2}{2}>5,
\]
because $R\ge37N/64$.

For $r=1$, $K_t^{(R)}(1)=R-2t$ and
$M_1=2\E(X-R)_+$ imply
\[
 u_{0,1}=2\E T-R=k-M_1>0.
\]
Here $k\ge1$ follows from $\alpha>2/5$, and $M_1<1/2$ follows from
\eqref{eq:half-norm}. Consequently,
\[
 R-|u_{0,1}|=R-k+M_1\ge n-2k
 \ge\frac{5N}{32}+1>5.
\]

Finally, $K_t^{(R)}(R-1)=(-1)^t(R-2t)$ yields
\[
 R-|u_{0,R-1}|\ge2\E\min\{T,R-T\}.
\]
On $\mathcal A$,
\[
 \min\{T,R-T\}\ge
 \min\left\{\frac n3,R-\frac{n-1}{2}\right\}
 \ge\frac{5N}{64},
\]
using \eqref{eq:kR-bounds}. Together with \eqref{eq:A-prob}, this gives
$R-|u_{0,R-1}|>5N/128\ge5$. This also proves the strict reserve at
$N=128$.
\end{proof}

\section{All dyadic blocks above the threshold}

\begin{theorem}[Shifted block theorem]\label{thm:large-block}
For every dyadic $N\ge128$, there are admissible integer packets satisfying
\eqref{eq:target-identity} for the shifted profile
\eqref{eq:critical-profile}. The resulting deterministic rule on
$[N,2N-1]$ is fair for every bias and satisfies \eqref{eq:main-profile}.
\end{theorem}
\begin{proof}
Take the real center of Proposition~\ref{prop:real-center}. At level zero
its norm is one, and every positive active level has norm less than
$1/2$ by Proposition~\ref{prop:digit-norms}. Hence the center satisfies
the packet boxes. By Lemma~\ref{lem:profile}, all structural assumptions
of Lemma~\ref{lem:lattice-rounding} hold, and the center vanishes on the
final two levels. Apply that rounding lemma.

Every packet boundary becomes $\pm1$ and is admissible. At level zero,
Lemma~\ref{lem:level-zero-reserve} supplies more than five units of
reserve at each nonboundary site. For $1\le i\le h$ and $0<r<R_i$,
\[
 \binom{R_i}{r}-|u_{i,r}|
 >\frac12\binom{R_i}{r}\ge\frac{R_i}{2}
 \ge\frac{37N}{128}>5.
\]
Thus the five-unit rounding error remains inside the packet boxes.
At levels $i>h$, the center is zero and the two consecutive suffixes
are zero. The sharper rounding bound gives $|f_{i,r}|\le2$ at interior
sites, whose capacities are at least two. The boundary sites are already
admissible. Packet realization proves fairness and
\eqref{eq:profile-stop} gives the deadline.
\end{proof}

\section{The six finite blocks}\label{sec:finite}

The analytic proof starts at $N=128$. The remaining blocks are
$N=2,4,8,16,32,64$, for which $s_N=0$.

\begin{proposition}[Exact finite certificate]\label{prop:finite-certificate}
For every dyadic $2\le N<128$, the profile \eqref{eq:critical-profile}
admits packets satisfying \eqref{eq:target-identity} and
\eqref{eq:packet-admissibility}.
\end{proposition}
\begin{proof}
The ancillary program \texttt{verify\_glazer\_critical.py} constructs these
six finite witnesses. Its SHA-256 digest is
\begin{center}
\small\texttt{07801d62aedf91b87e122bd96f53a9bd450d64dd4e265caaff364f3da0dced61}.
\end{center}
The packet arrays are also distributed separately in
\texttt{certificates/finite\_blocks.json}; their checker does not use the
folding or rounding construction.

Every mathematical comparison uses integers or rational numbers. First,
using rational endpoints whose squares straddle $5$ and the identity
\[
 \log x=2\sum_{j\ge0}\frac{z^{2j+1}}{2j+1},
 \qquad z=\frac{x-1}{x+1},
\]
with the remainder bound
\[
 0<2\sum_{j>M}\frac{z^{2j+1}}{2j+1}
 \le\frac{2z^{2M+3}}{(2M+3)(1-z^2)},
\]
the program certifies
\begin{equation}\label{eq:beta-enclosure}
 0.5979874356654401<\beta<0.5979874356654403.
\end{equation}
As an independent and purely algebraic certification of the floors,
the program also checks
\begin{equation}\label{eq:algebraic-floor}
 5^r<\phi^{2m}<5^{r+1},\qquad r=\lfloor\beta m\rfloor.
\end{equation}
These comparisons are equivalent to $r<\beta m<r+1$ because $\log5>0$.
Using $\phi^{2m}=F_{2m}\phi+F_{2m-1}$, they reduce to exact integercomparisons in $\mathbb Q(\sqrt5)$, followed by squaring only positive
quantities. Finally,
$\lfloor\alpha m\rfloor=m-1-\lfloor\beta m\rfloor$.
The finite floor witnesses are included in the JSON certificate.
Thus the decimal enclosure is not a necessary oracle for the finite
floor claims.
The program then uses \eqref{eq:Fnext-coeff} to construct the folded
center with denominator $2^{N-1}$, performs integral kernel rounding,
and checks every packet's capacity and parity. Finally, it expands the
left side of \eqref{eq:target-identity} coefficient by coefficient over
$\mathbb Z[z]$ and obtains exactly $(1,0,\ldots,0)$. It also checks
\eqref{eq:profile-stop} at every level. The independently checked packet
arrays are therefore witnesses to the proposition.

\begin{table}[ht]
\centering
\begin{tabular}{@{}rrrrr@{}}
\toprule
$N$ & $h$ & packet sites & upper bound on $\max|f-u|$ & least interior margin\\
\midrule
$2$  & $0$  & $3$    & $1$ & --\\
$4$  & $1$  & $9$    & $1$ & $2$\\
$8$  & $2$  & $32$   & $3$ & $2$\\
$16$ & $4$  & $122$  & $4$ & $2$\\
$32$ & $8$  & $474$  & $4$ & $2$\\
$64$ & $16$ & $1869$ & $5$ & $2$\\
\bottomrule
\end{tabular}
\caption{Exact finite witnesses. The interior margin is
$\binom{R_i}{r}-|f_{i,r}|$ for $0<r<R_i$; the error bounds are integers,
not rounded numerical tests.}
\label{tab:finite}
\end{table}
\end{proof}

\section{Proof of the main theorem}

\begin{proof}[Proof of Theorem~\ref{thm:main}]
For $L=1$, use $01\mapsto H$ and $10\mapsto T$. For $L\ge2$, let $N$
be its unique dyadic scale, so $N\le L<2N$, and write $L=N+i$.
Use Theorem~\ref{thm:large-block} if $N\ge128$, and
Proposition~\ref{prop:finite-certificate} otherwise.
Fix a definite admissible signing for each block, using the
lexicographic packet realization in Lemma~\ref{lem:packet-realization}.
Read through the first change, determine $N$ and $i$, then read $R_i$
additional bits and output the assigned sign. The resulting rule is
deterministic, uses no knowledge of $p$, and stops at $2N-a_i$.

The blocks partition the nonconstant sequences, and each block has zero
signed probability for every bias. Since $p,1-p<1$, the two infinite
constant sequences have probability zero. Countable additivity now gives
$\Prb_p(H)=\Prb_p(T)=1/2$. The deadline is \eqref{eq:profile-stop}.
Since $s_N\ge0$, this implies \eqref{eq:main-bound}.

Furthermore $s_N\ge\log_5N-4$ and $N>L/2$, so for $L\ge2$,
\[
 T(L)\le\lceil C_*L\rceil-s_N
 <C_*L-\log_5L+5+\log_52
 <C_*L-\log_5L+6.
\]
The same final bound holds directly for $L=1$. This proves
\eqref{eq:main-log-bound}.
\end{proof}

\section{Optimality without complement symmetry}\label{sec:optimality}

A rule is \emph{separately balanced} on a block if its two outputs have
equal probability on the event defining that block for every $p\in(0,1)$.
At horizon $H$, pad every decision to an $H$-letter word. This does not
require the rule to read the padding. A per-level deadline $t_i$ means
that the sign on the level $L=N+i$ is determined by the first $t_i$ bits.
Adaptive rules are included: take $t_i$ to be an upper bound on all their
stopping times at that level.

\begin{theorem}[Evaluation obstruction for arbitrary signings]\label{thm:obstruction}
Consider a deterministic, separately balanced block $[N,2N-1]$ at common
horizon $2N$. Suppose its per-level deadlines satisfy
$N+i+1\le t_i\le2N$. Set
\[
 a_i=2N-t_i,\qquad R_i=t_i-(N+i)-1.
\]
No complement symmetry is assumed. For every $0<\theta<1$,
\begin{equation}\label{eq:evaluation-obstruction}
 \sum_{i=0}^{N-1}\theta^i(1-\theta)^{a_i}
                  (1+\theta)^{N-1-i-a_i}\ge1.
\end{equation}
\end{theorem}
\begin{proof}
Use only the branch whose initial run consists of zeros. Let $e^+_{i,r}$
be its signed number of read tails of length $R_i$ with $r$ zeros, and put
\[
 E_i^+(x)=\sum_{r=0}^{R_i}e^+_{i,r}x^r,\qquad
 P_+(x)=\sum_{i=0}^{N-1}x^i(1+x)^{a_i}E_i^+(x).
\]
The coefficient of $x^\ell$ is the signed number of zero-run words with
$N+\ell$ zeros at horizon $2N$. For $\ell\ge1$, a word whose initial
run consists of ones cannot contribute: such a word has at most $N$
zeros. Lemma~\ref{lem:type-balance} therefore forces every nonconstant
coefficient of $P_+$ to vanish.

Its constant coefficient is the sign of the unique zero-run word with
exactly $N$ zeros, namely $0^N1^N$. Hence $P_+\equiv\pm1$.
Since $|e^+_{i,r}|\le\binom{R_i}{r}$, evaluation at $x=-\theta$ gives
\[
 1\le\sum_i\theta^i(1-\theta)^{a_i}|E_i^+(-\theta)|
 \le\sum_i\theta^i(1-\theta)^{a_i}(1+\theta)^{R_i}.
\]
This is \eqref{eq:evaluation-obstruction}.
\end{proof}

Put
\begin{equation}\label{eq:theta-star}
 \theta_*:=\phi^{-2}=\frac{3-\sqrt5}{2}.
\end{equation}
The identities
\[
 1-\theta_*=\phi^{-1},\qquad
 1+\theta_*=\frac{\sqrt5}{\phi},\qquad
 (\sqrt5)^\beta=\phi
\]
give an exact version of the obstruction in terms of the deadlines.

\begin{corollary}[Exact deadline inequality]\label{cor:deadline}
Under the assumptions of Theorem~\ref{thm:obstruction},
\begin{equation}\label{eq:deadline-inequality}
 \boxed{\displaystyle
 \sum_{i=0}^{N-1}(\sqrt5)^{\,t_i-C_*(N+i)}\ge1+\theta_* .}
\end{equation}
\end{corollary}
\begin{proof}
For $t_i=2N-a_i$, direct substitution yields
\[
 \theta_*^i(1-\theta_*)^{a_i}(1+\theta_*)^{N-1-i-a_i}
 =\frac{(\sqrt5)^{t_i-C_*(N+i)}}{1+\theta_*}.
\]
Multiply \eqref{eq:evaluation-obstruction} at $\theta=\theta_*$ by
$1+\theta_*$.
\end{proof}

\begin{corollary}[Sharp leading coefficient]\label{cor:framework-optimal}
Suppose fixed real constants $C,D$ satisfy $T(L)\le CL+D$ for a
universally fair extractor whose sufficiently large dyadic blocks are
separately balanced and have horizon $2N$. Then $C\ge C_*$. Thus the
infimum of affine coefficients in this class is $C_*$, and it is
attained, for example, with $D=1$.
\end{corollary}
\begin{proof}
Take per-level deadlines $t_i\le C(N+i)+D$. They are at least $N+i+1$
because no constant prefix can be terminal. If $C<C_*$, the left side
of \eqref{eq:deadline-inequality} is at most
\[
 N(\sqrt5)^{(C-C_*)N+D}\longrightarrow0,
\]
contradicting its fixed positive lower bound. The construction in
Theorem~\ref{thm:main} supplies the matching upper bound.
\end{proof}

\begin{corollary}[Restriction on second-order savings]\label{cor:second-order}
Every block in Theorem~\ref{thm:obstruction} satisfies
\begin{equation}\label{eq:log-obstruction}
 \max_{0\le i<N}\bigl(t_i-C_*(N+i)\bigr)
 \ge-2\log_5N+\log_{\sqrt5}(1+\theta_*).
\end{equation}
In particular, a uniform bound
\[
 T(L)\le C_*L-\lambda\log_5L+D
\]
within the separately balanced horizon-$2N$ class requires $\lambda\le2$.
If $\Lambda$ is the supremum of all coefficients $\lambda$ attainable
in this sense with some finite $D$, then
\begin{equation}\label{eq:lambda-range}
 1\le\Lambda\le2.
\end{equation}
\end{corollary}
\begin{proof}
If the maximum in \eqref{eq:log-obstruction} is $b$, then
\eqref{eq:deadline-inequality} gives $N(\sqrt5)^b\ge1+\theta_*$;
taking logarithms proves the first assertion. If $\lambda>2$, then
$N\le N+i$ gives
\[
 \sum_i(\sqrt5)^{t_i-C_*(N+i)}
 \le (\sqrt5)^D N^{1-\lambda/2}\longrightarrow0,
\]
a contradiction. Finally, \eqref{eq:main-log-bound} attains $\lambda=1$.
The supremum $\Lambda$ is not asserted to be attained.
\end{proof}

\section{Sublinear additional horizons}\label{sec:extra-horizon}

We next allow a common horizon $H_N=2N+j_N$ with nonnegative integers $j_N$. The argument
requires separate block balance, but still imposes no complement symmetry.

\begin{lemma}[Degree bound from separated run orientations]\label{lem:degree-window}
Let $N\ge1$ and $j\ge0$ be integers. For a separately balanced block
$[N,2N-1]$ at horizon $H=2N+j$, let
$t_i\in[N+i+1,H]$ be per-level deadlines. Put
$a_i=H-t_i$, $R_i=t_i-N-i-1$, and form the zero-run polynomial
\[
 P_+(x)=\sum_{i=0}^{N-1}x^i(1+x)^{a_i}E_i^+(x)
\]
as in Theorem~\ref{thm:obstruction}. Then
\[
 \deg P_+\le j,\qquad P_+(0)\in\{\pm1\}.
\]
\end{lemma}
\begin{proof}
The coefficient of $x^\ell$ is the signed count in the zero-run branch
with $N+\ell$ zeros. A one-run word has at most $H-N=N+j$ zeros.
Thus every coefficient with $\ell>j$ vanishes by type balance.
The constant coefficient is the sign of the unique word
$0^N1^{H-N}$.
\end{proof}

\begin{lemma}[Elementary polynomial extrapolation]\label{lem:extrapolation}
Let $0<a<b$, let $j\ge0$ be an integer, and let a real polynomial $P$
have degree at most $j$.
Then
\begin{equation}\label{eq:extrapolation}
 |P(0)|\le\left(\frac{6b}{b-a}\right)^j
          \max_{a\le\theta\le b}|P(-\theta)|.
\end{equation}
\end{lemma}
\begin{proof}
For $j=0$ the assertion is immediate. For $j\ge1$, interpolate at
$-\theta_k$, where $\theta_k=a+k(b-a)/j$ and $0\le k\le j$.
The sum of the absolute values at zero of the Lagrange basis polynomials
is at most
\[
 \sum_{k=0}^j\frac{b^j}{((b-a)/j)^j k!(j-k)!}
 =\left(\frac{2bj}{b-a}\right)^j\frac1{j!}.
\]
The elementary inequality $j^j/j!\le3^j$ completes the proof. To verify
that inequality without an asymptotic factorial estimate, note that
\[
 \frac{(m+1)^{m+1}/(m+1)!}{m^m/m!}
 =\left(1+\frac1m\right)^m
 \le\sum_{r=0}^m\frac1{r!}<3,
\]
using $r!\ge2^{r-1}$ for $r\ge1$; now induct from $m=1$.
\end{proof}

\begin{theorem}[Robust leading-coefficient obstruction]\label{thm:sublinear-horizon}
Suppose a universally fair deterministic extractor separately balances
all sufficiently large dyadic blocks and has common horizons
$H_N=2N+j_N$, where $j_N/N\to0$ along the powers of two. If
$T(L)\le CL+D$ for fixed real $C,D$, then $C\ge C_*$.
\end{theorem}
\begin{proof}
Assume $C<C_*$. With the notation of Lemma~\ref{lem:degree-window},
capacity gives, for $0<\theta<1$,
\begin{equation}\label{eq:extra-evaluation}
 |P_+(-\theta)|\le
 \sum_{i=0}^{N-1}\theta^i(1-\theta)^{H_N-t_i}
                     (1+\theta)^{t_i-N-i-1}.
\end{equation}
Each summand is increasing as a function of $t_i$, because
$(1+\theta)/(1-\theta)>1$. Replacing $t_i$ by its upper bound
$C(N+i)+D$ and writing $v=i/N$ bounds the summand by
\[
 \frac1{1+\theta}
 \left(\frac{1+\theta}{1-\theta}\right)^D
 (1-\theta)^{j_N}\exp\{N G_C(\theta,v)\},
\]
where
\begin{align*}
 G_C(\theta,v)
 &=v\log\theta+[2-C(1+v)]\log(1-\theta)\\
 &\hspace{2em}+[C(1+v)-1-v]\log(1+\theta).
\end{align*}
At the critical evaluation point,
\[
 G_C(\theta_*,v)=(C-C_*)(1+v)\log\sqrt5<0
 \qquad(0\le v\le1).
\]
Because $G_C$ is affine in $v$, it suffices to check its two endpoint
values. By continuity, there are $0<a<\theta_*<b<1$ and $\delta>0$
for which $G_C(\theta,v)\le-\delta$ whenever $a\le\theta\le b$ and
$0\le v\le1$. The remaining prefactor is bounded on $[a,b]$, and
$(1-\theta)^{j_N}\le1$. Thus, with a constant $K$ independent of $N$,
\[
 \max_{a\le\theta\le b}|P_+(-\theta)|\le KN e^{-\delta N}.
\]
Lemma~\ref{lem:degree-window} and Lemma~\ref{lem:extrapolation} imply
\[
 1=|P_+(0)|\le
 \left(\frac{6b}{b-a}\right)^{j_N}KN e^{-\delta N}.
\]
The logarithm of the right side tends to $-\infty$, since
$j_N=o(N)$ and $\log N=o(N)$. This is the required contradiction.
\end{proof}

\begin{remark}[Scope of the obstruction]
Dropping complement symmetry does not circumvent these bounds, nor does
an $o(N)$ increase of the common block horizon. The arguments do not
exclude cancellation between different dyadic blocks or separately
balanced rules with larger additional horizons. In particular, they do
not prove that $C_*$ is optimal among unrestricted global extractors.
\end{remark}

\section{Exchangeable sources and conditional fairness}\label{sec:exchangeability}

The finite combinatorial balance is stronger than fairness for independent
flips. It is useful to state that consequence without any mixture
representation theorem.

\begin{theorem}[Exchangeable-source extension]\label{thm:exchangeable}
The deterministic extractor of Theorem~\ref{thm:main} is fair under every
law on infinite binary sequences whose finite-dimensional distributions
are invariant under coordinate permutations, provided $L<\infty$ almost
surely. Its pathwise deadlines are unchanged.
Moreover, for a dyadic $N\ge2$, let $K_N$ be the number of zeros among
the first $2N$ inputs. Whenever the conditioning event has positive
probability,
\begin{equation}\label{eq:conditional-type-fairness}
 \Prb(H\mid N\le L<2N,\ K_N=k)=\frac12.
\end{equation}
The corresponding statement for $N=1$ follows from the pair $01,10$.
\end{theorem}
\begin{proof}
By construction, each block signing has zero signed probability for all
i.i.d. biases. Lemma~\ref{lem:type-balance} says that its signed count
vanishes separately in every Hamming-weight class at horizon $2N$.
Under an exchangeable finite marginal, all words in a fixed Hamming-weight
class have the same probability, because a coordinate permutation maps
any one of them to any other. Hence the two output events have equal
probability within the intersection of the block and that class. This
proves \eqref{eq:conditional-type-fairness}.
Summing over the weight classes gives fairness on each block, and summing
over blocks gives total fairness because $L<\infty$ almost surely.
The stopping statement is pathwise and does not depend on the source law.
\end{proof}

\begin{example}
Choose a random parameter $P$ supported in $(0,1)$, then conditionally
on $P$ generate i.i.d. inputs of bias $P$. The resulting source is
exchangeable and has $L<\infty$ almost surely, so the theorem applies
although the inputs need not be independent after averaging over $P$.
\end{example}

\begin{remark}
The condition on constant sequences is necessary for almost-sure
termination of this rule. Also, \eqref{eq:conditional-type-fairness} does
not assert fairness conditional on each individual value of $L$, nor
independence of the output from the stopping time.
\end{remark}

\section{Why dyadic blocks occur at horizon \texorpdfstring{$2N$}{2N}}\label{sec:dyadic-necessity}

Here $N$ may be any positive integer, not necessarily a power of two.

\begin{theorem}[Dyadic necessity and sufficiency]\label{thm:dyadic-iff}
The set of $2N$-letter words whose initial run length lies in $[N,2N-1]$
admits a deterministic signing that is balanced for every bias if and
only if $N$ is a power of two.
\end{theorem}
\begin{proof}
For $0\le\ell<N$, the number of zero-run words with $N+\ell$ zeros is
\[
 \sum_{i=0}^{N-1}\binom{N-1-i}{\ell-i}=\binom N\ell,
\]
where out-of-range binomial coefficients are zero. This follows by
extracting coefficients from
\[
 \sum_{i=0}^{N-1}x^i(1+x)^{N-1-i}=(1+x)^N-x^N.
\]
For $1\le\ell<N$, these are all the words with that zero count; the
one-run branch contributes only to lower counts. The reflected classes
have the same cardinalities. The central class contains exactly the two
words $0^N1^N$ and $1^N0^N$.

By Lemma~\ref{lem:type-balance}, a fair deterministic split exists
exactly when every intermediate coefficient $\binom N\ell$ is even.
If $N$ is a power of two, repeated squaring over $\mathbb F_2$ gives
$(1+x)^N=1+x^N$, so all such coefficients are even. Conversely, writing
$N=\sum_j\epsilon_j2^j$ in binary gives
\[
 (1+x)^N=\prod_{j:\epsilon_j=1}(1+x^{2^j})\quad\text{in }\mathbb F_2[x].
\]
If two or more binary digits are nonzero, the factor with the least
exponent supplies a nonzero coefficient strictly between degrees zero
and $N$. Thus $N$ must be a power of two. When it is, split every finite
Hamming-weight class into two equal parts; reading the full horizon
realizes such a signing.
\end{proof}

\begin{corollary}[No fixed two-bit output in one such block]\label{cor:no-two-bit}
At common horizon $2N$, no deterministic rule on $[N,2N-1]$ can have four
outputs with equal conditional probabilities within the block for every
bias. In particular, the one-bit construction cannot be replaced by two
uniform output bits per block without changing these requirements.
\end{corollary}
\begin{proof}
Equal output probabilities for every bias force equal counts in each
Hamming-weight class, by applying the polynomial argument to each pair
of outputs. The central class has two words and cannot be divided into
four equal integer parts.
\end{proof}

\section*{Reproducibility and computational scope}
\addcontentsline{toc}{section}{Reproducibility and computational scope}

The package contains this source, the compiled PDF,
\texttt{verify\_glazer\_critical.py}, the finite packet arrays in
\texttt{certificates/finite\_blocks.json}, and the recorded verification
output. The six arrays with $N<128$ are the finite witnesses used in
Section~\ref{sec:finite}. Larger test blocks are supplementary checks;
the all-$N$ theorem is the mathematical argument of
Sections~\ref{sec:profile}--\ref{sec:finite}, not an extrapolation from
those tests.

Run the required finite checks and their independent witness checker by
\begin{verbatim}
python scripts/verify_glazer_critical.py
python scripts/verify_glazer_critical.py \
  --check-certificate scripts/certificates/finite_blocks.json
\end{verbatim}
The recorded extended run is reproduced by
\begin{verbatim}
python scripts/verify_glazer_critical.py --extra 128 256 512 1024 \
  --export scripts/certificates/finite_blocks.json
\end{verbatim}
Do not pass Python's \texttt{-O} flag: assertions are part of the checker.
All mathematical comparisons use exact integer or rational arithmetic;
decimal displays are produced only after the corresponding exact checks.
A numerical certificate is not, by itself, a proof checked by a
proof-assistant kernel.

\paragraph{Character expansion.}
For $c_r=[x^r](1-x)^t(1+x)^{R-t}$, the verifier uses
\[
 c_0=1,\quad c_{-1}=0,\quad
 (r+1)c_{r+1}=(R-2t)c_r-(R-r+1)c_{r-1}.
\]
This follows by differentiating the product and clearing $1-x^2$.
The script checks the recurrence independently against direct polynomial
multiplication for $0\le t\le R\le40$.

\paragraph{Exact large-block polynomial checking.}
The required finite certificates are checked coefficient by coefficient.
For supplementary blocks above $128$, the script may instead use the
following exact test. If an integer polynomial $Q$ has each coefficient
bounded in magnitude by $M$ and $B>M+1$ is an integer, then
\[
 Q(B)=0\quad\Longrightarrow\quad Q=0.
\]
Indeed, a nonzero leading integer coefficient has magnitude at least
one, and its contribution dominates all lower terms, whose total
magnitude is at most $M(B^d-1)/(B-1)<B^d$. For the difference from the
target packet identity, the script uses the explicit bound
\[
 M=|\text{target}|+\sum_i2^{a_i}\sum_r|c_{i,r}|.
\]
It clears the known denominator before checking a rational center.
This is an injective integer evaluation, not a probabilistic fingerprint
or a congruence modulo a chosen prime.

\begin{thebibliography}{99}

\bibitem{GlazerPuzzling2024}
E.~Glazer,
\newblock Extracting a fair coin flip from a biased die, with a time limit,
\newblock \emph{Puzzling Stack Exchange}, August 8, 2024,
\newblock \href{https://puzzling.stackexchange.com/questions/127763/extracting-a-fair-coin-flip-from-a-biased-die-with-a-time-limit}
{puzzling.stackexchange.com/questions/127763}.

\bibitem{WangPuzzling2024}
C.~Wang,
\newblock Answer to Glazer's ``Extracting a fair coin flip from a biased die, with a time limit,''
\newblock \emph{Puzzling Stack Exchange}, August 8, 2024,
\newblock \href{https://puzzling.stackexchange.com/questions/127763/extracting-a-fair-coin-flip-from-a-biased-die-with-a-time-limit}
{puzzling.stackexchange.com/questions/127763}.

\bibitem{WagonPuzzling2025}
S.~Wagon,
\newblock Answer to Glazer's ``Extracting a fair coin flip from a biased die, with a time limit,''
\newblock \emph{Puzzling Stack Exchange}, April 26, 2025,
\newblock with a formulaic construction attributed there to S.~Wagon and P.~Winkler,
\newblock \href{https://puzzling.stackexchange.com/questions/127763/extracting-a-fair-coin-flip-from-a-biased-die-with-a-time-limit}
{puzzling.stackexchange.com/questions/127763}.

\bibitem{GlazerAMM2026}
E.~Glazer,
\newblock Problem 12592,
\newblock in D.~H. Ullman, D.~J. Velleman, S.~Wagon, and D.~B. West,
``Problems and Solutions,''
\emph{Amer. Math. Monthly} \textbf{133} (2026), no.~3, 290--300, p.~291,
\newblock \href{https://doi.org/10.1080/00029890.2026.2599721}
{doi:10.1080/00029890.2026.2599721}.

\bibitem{HoeffdingSimons1970}
W.~Hoeffding and G.~Simons,
\newblock Unbiased coin tossing with a biased coin,
\newblock \emph{Ann. Math. Statist.} \textbf{41} (1970), no.~2, 341--352,
\newblock \href{https://doi.org/10.1214/aoms/1177697074}
{doi:10.1214/aoms/1177697074}.

\bibitem{Elias1972}
P.~Elias,
\newblock The efficient construction of an unbiased random sequence,
\newblock \emph{Ann. Math. Statist.} \textbf{43} (1972), no.~3, 865--870,
\newblock \href{https://doi.org/10.1214/aoms/1177692552}
{doi:10.1214/aoms/1177692552}.

\bibitem{StoutWarren1984}
Q.~F. Stout and B.~Warren,
\newblock Tree algorithms for unbiased coin tossing with a biased coin,
\newblock \emph{Ann. Probab.} \textbf{12} (1984), no.~1, 212--222,
\newblock \href{https://doi.org/10.1214/aop/1176993384}
{doi:10.1214/aop/1176993384}.

\bibitem{Peres1992}
Y.~Peres,
\newblock Iterating von Neumann's procedure for extracting random bits,
\newblock \emph{Ann. Statist.} \textbf{20} (1992), no.~1, 590--597,
\newblock \href{https://doi.org/10.1214/aos/1176348543}
{doi:10.1214/aos/1176348543}.

\bibitem{JuelsEtAl2000}
A.~Juels, M.~Jakobsson, E.~A.~M. Shriver, and B.~K. Hillyer,
\newblock How to turn loaded dice into fair coins,
\newblock \emph{IEEE Trans. Inform. Theory} \textbf{46} (2000), no.~3, 911--921,
\newblock \href{https://doi.org/10.1109/18.841170}
{doi:10.1109/18.841170}.

\bibitem{PaeLoui2005}
S.-I.~Pae and M.~C. Loui,
\newblock Optimal random number generation from a biased coin,
\newblock in \emph{Proceedings of the Sixteenth Annual ACM--SIAM Symposium
on Discrete Algorithms (SODA 2005)}, 1079--1088,
\newblock \href{https://doi.org/10.1145/1070432.1070587}
{doi:10.1145/1070432.1070587}.

\bibitem{vonNeumann1951}
J.~von Neumann,
\newblock Various techniques used in connection with random digits,
\newblock in A.~S. Householder, G.~E. Forsythe, and H.~H. Germond (eds.),
\emph{Monte Carlo Method}, National Bureau of Standards Applied Mathematics
Series 12, U.S. Government Printing Office, Washington, DC, 1951,
pp.~36--38.

\end{thebibliography}

\end{document}